\documentclass[11pt,a4paper]{article}
\usepackage[utf8]{inputenc}
\usepackage[T1]{fontenc}
\usepackage{lmodern,amsmath,amssymb,textcomp}
\usepackage[margin=24mm]{geometry}
\usepackage{longtable,booktabs,array,enumitem}
\usepackage[hidelinks]{hyperref}
\setlength{\parindent}{0pt}
\setlength{\parskip}{6pt}
\emergencystretch=3em
\begin{document}
{\LARGE\bfseries A log(5/3) uniform bound for representations as two nonnegative cubes\par}\medskip
{\small Provisional mathematical authors: Rachel Teo, Wenyi Wang, Hamdi Abderrahmene\par}\medskip
{\small Judging and evaluation record: the submissions discussed in this document have been judged and evaluated by Alejandro Zarzuelo Urdiales for OpenMath 2026. This dated review records the stated verification, classification and unresolved holds; it is a preliminary evaluation rather than a signed award decision.\par}

{\small Event provenance: OpenMath 2026 ran 27 September-2 October 2026 with worldwide online participation. Detailed organizer listings specify Harvard Innovation Labs for the 27 September in-person event; the OpenMath programme advertises a hybrid kickoff at MIT. Sundai identifies its Harvard/MIT community. Organizer sources: \url{https://rsihouse.ai/openmath} ; \url{https://luma.com/yzp9abvr} ; \url{https://www.sundai.club/events/boston/recursive-learning-hack-with-harvard-innovation-labs}\par}

{\small Rachel Teo  -  Sakana AI.\par}

{\small Wenyi Wang  -  Sakana AI.\par}

{\small Hamdi Abderrahmene  -  Sakana AI.\par}

{\small Affiliations follow the recorded author declarations or public institutional/profile sources. Preferred author names, author order and publication affiliations await author review. This editorial compilation does not establish anthology coauthorship.\par}

{\small Original-result working manuscript | OpenMath 2026 | 5 October 2026\par}\medskip
{\small Central source and adjudication archive: \url{https://github.com/alejandrozu/openmath-2026-judging}\par}

\subsection{Abstract}
For ordered nonnegative representations x\ensuremath{^3}+y\ensuremath{^3}=n, the formal endpoint gives the eventual bound exp((log(5/3)+\ensuremath{\varepsilon})log n/log log n). The entrant-supplied ordinary proof optimizes the same weighted ternary coding method to coefficient C=log 3\ensuremath{-}H((2\ensuremath{-}\ensuremath{\surd}2)/3)\ensuremath{-}((2\ensuremath{-}\ensuremath{\surd}2)/3)log 2\ensuremath{\approx}0.4695030882. A gcd transfer and prime cutoff remove primitive and cube-free restrictions. The stronger ordinary endpoint is not asserted as Lean-checked.

\subsection{Count and uniform theorem}
Let R(n)=\#\{(x,y)\ensuremath{\in}\ensuremath{\mathbb{N}}\ensuremath{^2}:x\ensuremath{^3}+y\ensuremath{^3}=n\}, with ordered pairs and zero coordinates included. For every \ensuremath{\varepsilon}>0 there is N such that R(n)\ensuremath{\leq}exp((log(5/3)+\ensuremath{\varepsilon})log n/log log n) for every n\ensuremath{\geq}N. No cubefree-target restriction or gcd-one restriction appears in the final theorem.

The coefficient log(5/3)\ensuremath{\approx}0.510826 improves the earlier same-family log(7/4)\ensuremath{\approx}0.559616 endpoint. The located approximately0.556477 comparison applies only to cubefree targets. It therefore cannot be substituted as a stronger all-target theorem. The result still falls far short of a fixed polylogarithmic bound, which remains the original target.

\subsection{Primitive ratios as a separated ternary code}
For a primitive positive solution, the ratio of its coordinates is a unit modulo the relevant prime powers, with cube equal to minus one. The unit cube-root lemma gives at most three choices at each prime-power coordinate, including the ramified prime three. Distinct solutions are encoded by distinct ternary words.

If two encoded words agree at selected prime-power coordinates, those moduli divide a cross determinant of the two solutions. An upper bound for the determinant and the target-size product force a weighted separation property. This converts an arithmetic representation problem into a weighted code problem without assuming the desired representation bound.

Elias/Plotkin localization bounds the size of a separated code in a weighted ball. An exact binomial block certificate uses blocks of 500 coordinates and support 90, producing the rational base 5/3 with an explicit finite multiplicative constant. The constant is harmless for the eventual leading exponent, but its proof is retained rather than replaced by a decimal numerical guess.

\subsection{From primitive pairs to the canonical count}
A general positive pair has gcd d and a primitive pair representing n/d\ensuremath{^3}. The source proves the exact gcd-fibre identity and sums over positive cube divisors d. Coordinatewise multiplicative weights bound this sum uniformly across targets, including those with repeated prime factors.

The prime cutoff and elementary log budgets control the resulting product. They reuse background modules from the k\ensuremath{\sigma}(k) development, while the cubic encoding and local-factor estimates are specific to this result. Finally at most two boundary pairs have a zero coordinate, and their contribution is absorbed. A semantic bridge identifies the resulting count with the original ordered cube-value representation definition; it does not import the open conjecture as a hypothesis.

\subsection{Verification, prior work and manuscript obligations}
The source contains 33 proof modules and two semantic/axiom audit modules in the recorded replay. Some arithmetic files were reconstructed after a workspace replacement and are clearly labelled variants; the dated fresh archived receipt supersedes their stale pending comments, without pretending the missing historical bytes were recovered. The final transitive axioms in those records are standard only.

Merikoski's ratio-root/determinant methods and classical coding/entropy bounds must be credited. The proposed new delta is the exact 5/3 formal weighted-code constant, the optimized ordinary coefficient 0.4695030882..., and their all-target transfer. This is a quantitative partial-bound paper, not a solution of Erdős 829 or a signed-integer representation theorem. A human exposition should display the weighted-distance condition, the finite block inequality and the gcd factorization so that the formal chain remains intelligible.

\subsection{Fresh independent verification of the displayed finite/formal scope}
Fresh independent replay completed 2026-10-05T01:22:59Z: all 31 frozen authored sources compiled, including the archived entrant-delivered reconstructed variants, under exact Lean 4.33.1 and Mathlib 0df444a360eaa60ab8c11dca51a86af692955474. Both actual requested axiom prints, canonicalCubeCount\_eventually\_exp and canonicalCubeCount\_eventual\_threshold, are present and contain only propext, Classical.choice and Quot.sound, with no admitted, native-evaluation or unknown axioms. This verifies the displayed all-target log(5/3) endpoint; the sharper ordinary coefficient .46950308823856934 remains a separately explained source result. The replay proves the delivered reconstructed proof, without claiming recovery of the first lost source bytes. Receipt: Verification/sakana-erdos829-log-five-thirds-fresh-build.json.

{\small This short manuscript records the construction and selected endpoints. The companion Original results anthology contains the extended ordinary proof exposition and its explicitly stated premises; the preserved Lean source remains the complete formal proof record.\par}

\subsection{Verification and reproducibility}
Fresh execution: sakana-erdos829-log-five-thirds: PASS; 31/31 custom modules compiled successfully; receipt Verification/sakana-erdos829-log-five-thirds-fresh-build.json. Compiler pins, source hashes and endpoint axiom outputs are in Verification. Historical receipts and finite checks do not replace fresh replay.

\begin{samepage}
\subsection{Erdos829FiveThirds.canonicalCubeCount\_eventually\_exp}
\begin{quote}\small\ttfamily theorem canonicalCubeCount\_eventually\_exp \{\ensuremath{\varepsilon} : \ensuremath{\mathbb{R}}\} (h\ensuremath{\varepsilon} : 0<\ensuremath{\varepsilon}) :\newline     \ensuremath{\forall}\ensuremath{^{f}} n : \ensuremath{\mathbb{N}} in atTop, (canonicalCubeCount n : \ensuremath{\mathbb{R}}) \ensuremath{\leq}\newline       Real.exp ((Real.log (5/3 : \ensuremath{\mathbb{R}})+\ensuremath{\varepsilon})*Real.log (n : \ensuremath{\mathbb{R}})/Real.log (Real.log (n : \ensuremath{\mathbb{R}})))\end{quote}
{\small Source: proofs/erdos829-log-five-thirds/FiveThirdsCanonical.lean:8.\par}

{\small Source SHA-256: 51541747f8104588463d8b7b051d69692876cee605c68a2f50df9da132e6a006\par}

\end{samepage}
\begin{samepage}
\subsection{Erdos829FiveThirds.canonicalCubeCount\_eventual\_threshold}
\begin{quote}\small\ttfamily theorem canonicalCubeCount\_eventual\_threshold \{\ensuremath{\varepsilon} : \ensuremath{\mathbb{R}}\} (h\ensuremath{\varepsilon} : 0<\ensuremath{\varepsilon}) :\newline     \ensuremath{\exists} N : \ensuremath{\mathbb{N}}, \ensuremath{\forall} n : \ensuremath{\mathbb{N}}, N\ensuremath{\leq}n \ensuremath{\to} (canonicalCubeCount n : \ensuremath{\mathbb{R}}) \ensuremath{\leq}\newline       Real.exp ((Real.log (5/3 : \ensuremath{\mathbb{R}})+\ensuremath{\varepsilon})*Real.log (n : \ensuremath{\mathbb{R}})/Real.log (Real.log (n : \ensuremath{\mathbb{R}})))\end{quote}
{\small Source: proofs/erdos829-log-five-thirds/FiveThirdsCanonical.lean:14.\par}

{\small Source SHA-256: 51541747f8104588463d8b7b051d69692876cee605c68a2f50df9da132e6a006\par}

{\small\textbf{Sources and attribution:} \href{https://www.erdosproblems.com/829}{1. www.erdosproblems.com / 829}; \href{https://github.com/google-deepmind/formal-conjectures/blob/main/FormalConjectures/ErdosProblems/829.lean}{2. google-deepmind/formal-conjectures / 829.lean}; \href{https://github.com/google-deepmind/formal-conjectures/blob/main/FormalConjecturesForMathlib/Combinatorics/Additive/Convolution.lean}{3. google-deepmind/formal-conjectures / Convolution.lean}; \href{https://arxiv.org/html/2112.03617}{4. arXiv source: 2112.03617}; \href{https://arxiv.org/html/2207.12487v2}{5. arXiv source: 2207.12487v2}; \href{https://github.com/alejandrozu/ulam-opdp-difficulty-atlas}{6. alejandrozu/ulam-opdp-difficulty-atlas}; \href{https://github.com/alejandrozu/openmath-2026-judging/blob/d0ed487f487fa7ec814ff705ab55249983fe8707/OpenMath-Judging/review/entries/E02-erdos829-log-five-thirds.txt}{7. alejandrozu/openmath-2026-judging / E02-erdos829-log-five-thirds.txt}; \href{https://github.com/alejandrozu/openmath-2026-judging}{Central source and adjudication archive}.\par}

\end{samepage}
\end{document}
